\section{Pre-Analysis Tab}
\label{sec:preanalysis}

The \app{Pre-Analysis} tab loads measured data, gates the
time-of-flight signal, bins the resulting spectrum on either
acceleration voltage or Doppler-shifted rest-frame frequency, and
overlays trial HFS models so the operator can read off initial
parameters before the full fit.

Each pre-analysis is a sub-tab under \app{+ Create Pre-Analysis}; the
sub-tab name (e.g.\ \app{Rh} in
Figure~\ref{fig:quickstart_preanalysis}) is shared with the Analysis
tab when settings are pushed across (Section~\ref{sec:analysis}). The
layout has two panes: a left column holding the \app{Data Files}
list, the \app{Plot Options} group, and the \app{HFS Models} panel,
and a wide centre pane with three plot sub-tabs --- \app{Spectrum},
\app{Calibrations}, and \app{Cooler Voltage}.

Three view controls sit in the corner of the sub-tab row and apply to
\emph{every} open pre-analysis project at once: a \app{Layout} combo
(\app{3 row stacked} = ToF / Spectrum / Timestamp as rows, or
\app{2 row stacked} = Spectrum and ToF side by side with the
Timestamp strip below, the original Data Viewer arrangement),
\app{X} / \app{Y} grid toggles, and a sun-moon button that switches
the plots to a \emph{dark mode} (black canvas, white chrome, neon
line colours). All three persist --- the dark state both in the
global settings and in the project save file.

\subsection{Loading data}

Click \app{Open\dots} to add one or more IGISOL-style ASDF run files;
\code{clstools}~\cite{clstools} parses each one and a row is added to
the \app{Data Files} list with a checkbox, an auto-assigned colour,
and a short header showing the cooler voltage $V$, laser setpoint
$\lambda$, and (after loading) the run duration. Tick the runs you
want plotted; the panels redraw immediately. \app{Remove checked}
drops the ticked rows, and the tri-state \app{Check all} box flips
every row at once. A circular-arrow button re-reads a still-growing
run from disk.

Default run colours cycle through the classic Windows~98 palette
(navy, maroon, green, purple, \dots), chosen for contrast between
consecutively loaded runs. Every file (and every HFS model) carries
\emph{two} colours --- one for the light canvas and one for dark
mode --- edited together in a dual-mode colour dialog (preset
palettes, hex/\code{R,G,B} entry, and saved custom slots shared
across the app). Clicking a row's colour swatch opens a small menu
with the colour dialog, an opacity slider, and a line-style selector.

Right-clicking a run row offers \app{Remove}, \app{Check}/\app{Uncheck},
\app{Set centroid offset\dots} (a manual MHz shift consumed by
frequency-domain merges), \app{Filter scans\dots}
(Section~\ref{sec:pa_scanfilter}), \app{Calibration\dots} and
\app{Calibration overview\dots}
(Section~\ref{sec:pa_calibration}). Merged entries instead offer
\app{View\dots}, \app{Edit\dots}, and \app{Export ASDF\dots}.

Different runs may sample different parts of the rest-frame spectrum
even when the lab-frame voltage range is identical, because each
run's cooler voltage and laser setpoint may differ. \denis{} keeps
each run's metadata so the rest-frame conversion is always applied
per run, never globally.

\subsection{The time-of-flight and timestamp gates}

The TOF panel shows the time-resolved ion signal for the selected
runs, histogrammed with an adjustable \app{Bin} size (default
1~\si{\micro\second}, down to 1~ns for narrow bunches). Enable the
\app{Gate} checkbox and set the window in \si{\micro\second} by
dragging the orange region on the plot or typing into the two value
boxes; only counts inside the window are histogrammed into the
spectrum below, which re-bins live as the gate moves.

The Timestamp panel plots event counts over wall-clock time with a
selectable \app{Unit} (Seconds / Minutes / Hours / Days) and its own
\app{Bin} size. Its \app{Gate} restricts the spectrum to a chosen
time window --- the TOF and timestamp histograms themselves always
show all events; only the spectrum is filtered.

The arrow-step convention is shared by the HFS-model value boxes.
Clicking inside a value box places the cursor on a digit, and the
up / down arrows then step the value by the order of magnitude of
that digit --- click a fine digit for small adjustments, a coarse one
for big jumps. The active digit is highlighted.

\subsection{Display axis and physics inputs}

The \app{Plot Options} group configures the display. The \app{X-axis}
drop-down offers five modes: \app{Voltage} (raw DAC scanning
voltage), \app{Calibrated voltage} (mapped through the calibration
table), \app{Calibrated beam energy} (cooler minus calibrated
scanning voltage), \app{Wavenumber}, and \app{Frequency} (rest-frame
MHz). The Doppler conversion uses the per-run cooler voltage, laser
setpoint, transition energies (\app{E lower} / \app{E upper}, with
derived \app{Transition} and \app{Fundamental} read-outs), the laser
\app{Harmonic}, and the isotope mass --- looked up from $Z$/$A$, with
an \app{Override} tick to type a custom amu value. When the
effective bin domain is frequency, the combo is restricted to
\app{Wavenumber} and \app{Frequency}: a frequency-bin centre has no
clean inverse back to volts.

\app{Channels} selects which PMTs (1--4, DC) to sum, \app{Normalize}
scales each run to unit peak height, and the \app{Cooler (V)} /
\app{Laser (cm$^{-1}$)} fields override the per-run metadata for all
files at once (a down-arrow button auto-fills from the selected
file). The override is useful as a sanity check and for runs whose
recorded values are known to be wrong.

\subsection{Binning}
\label{sec:pa_binning}

The spectrum is binned \emph{per scan step} by default: one bin per
native DAC voltage step (or its Doppler-shifted frequency
equivalent), which reproduces the legacy Data Viewer exactly and is
structurally immune to the uniform-grid aliasing that produces
doubled-count spikes when a fixed grid beats against the scan step
spacing. The \app{Bin:\ $N\times$ step} spin in the Spectrum header
groups $N$ adjacent steps into one bin (count-weighted centres,
summed counts) when more smoothing is wanted. Pre-Analysis shares the
same binning engine as the Analysis Source block
(Section~\ref{sec:analysis}), so a spectrum read here matches what
the fitter sees.

\subsection{HFS model overlay}

The \app{HFS Models} panel (left column) adds trial hyperfine
models. Each model is a checkable panel with a name, dual
light/dark colours, a line-style selector, an opacity slider, and
the parameter set: $I$, $J_l$, $J_u$ (one row), then $A_l$, $A_u$,
$B_l$, $B_u$, \app{Centroid}, \app{Scale}, \app{Bkg}, and Gaussian /
Lorentzian FWHMs. \app{Fix Al/Au} and \app{Fix Bl/Bu} lock the pairs
to a typed ratio --- handy when the upper-state constant is poorly
constrained; the locks carry over as fixed flags when the model is
imported into the Analysis tab.

Each parameter row has a value box, a slider (with editable limits
via the \app{\dots} button), and digit-aware arrow stepping. Tick
\app{Show Labels} to annotate each component on the spectrum with its
transition label, and \app{Peak Amplitudes} to expose a per-component
editor: each peak can stay on its \app{Racah} intensity, be typed
freely (\app{Free}), or be \app{Linked} to another peak by a ratio.

To overlay a ground state and an isomer, add two models and rename
them (e.g.\ \app{g.s.} and \app{meta} in
Figure~\ref{fig:quickstart_preanalysis}). \app{Duplicate checked}
copies ticked models; \app{Remove checked} removes them (undoable
with \code{Ctrl+Z}).

\subsection{Per-scan filtering}
\label{sec:pa_scanfilter}

A \emph{scan} is one full sweep of the voltage range. Right-click a
run and choose \app{Filter scans\dots} to drop individual low-quality
scans before binning: the dialog lists every scan with its start
time, duration, event count, PMT count and rate, with bulk
check/uncheck controls and an \app{Exclude rate $<$} threshold.
Alternatively tick \app{Show scans} on the Timestamp panel (single
non-merged run only), pan/zoom to a bad stretch, and click
\app{Exclude scans in view}. Exclusions live in a registry shared
with the Analysis tab --- a filter set here applies to the same file
there --- and persist under the save file's top-level
\code{scan\_filters} key.

\subsection{Virtual splits}

\menu{Edit $\rightarrow$ Split File\dots} turns one parent ASDF into
two virtual sub-runs described by tiny \code{.vasdf} sidecar files,
each carrying a raw-voltage gate plus its own cooler/laser metadata
--- useful when two isotopes were scanned in one run. Loading a
\code{.vasdf} in Pre-Analysis or Analysis behaves like an ordinary
run with a permanent voltage gate; the parent file is never
modified.

\subsection{Merging runs}

When several runs measure the same isotope, tick at least two loaded,
non-merged runs and click \app{Merge Checked}. The shared merge
dialog (Section~\ref{sec:analysis}) asks whether to merge in the
lab-frame \app{voltage} domain or the Doppler-shifted \app{frequency}
domain --- the default follows the current x-axis. Use voltage-domain
merges when the runs share cooler voltage and laser setpoint, and
frequency-domain merges otherwise (each file is shifted to the rest
frame with its own metadata first; per-file centroid offsets and any
GP centroid alignment are applied at the same stage). The result is a
synthetic entry that bins, overlays models, and exports like an
ordinary run.

\subsection{Calibrations sub-tab}

The \app{Calibrations} sub-tab shows three stacked diagnostics of the
scanning-voltage DAC, plotted against the requested set voltage for
every checked, non-split run: the readback response, the
readback$-$set offset (supply offset and drift), and the
step-to-step readback change (step uniformity). It is a quick health
check before pulling a run into a fit; the per-run calibration
\emph{fit} and its overrides live in the dialogs described in
Section~\ref{sec:pa_calibration}.

\subsection{Cooler Voltage sub-tab}

The \app{Cooler Voltage} sub-tab diagnoses the stability of the
cooler high voltage over each run using robust statistics (median
reference $V_{\mathrm{ref}}$ and $\sigma = 1.4826\cdot$MAD, immune to
spikes). A sortable per-run summary table reports
$V_{\mathrm{ref}}$, $\sigma$, peak-to-peak, spike count, signed
extremes, and $\sigma$ converted to its approximate Doppler impact in
MHz. Below it, three panes show the raw per-event cooler voltage,
the deviation from the run average (with $\pm\sigma$ / $\pm3\sigma$
bands, a per-bin median drift line, spike dots, and a secondary
$\approx\Delta\nu$~(MHz) axis), and the per-bin ripple amplitude
(RMS and P95$-$P5). \app{Clip y to $\pm4\sigma$} keeps normal ripple
readable when spikes are present; \app{Bins} sets the time binning.
Virtual splits are skipped on both diagnostic sub-tabs --- they share
their parent's arrays.

\subsection{Per-run voltage calibration}
\label{sec:pa_calibration}

Each run carries a DAC$\rightarrow$HV calibration table; a polynomial
through it turns every event's DAC value into a real voltage, so a
bad table biases the whole frequency axis (at 30~kV and mass 51,
1~V of calibration error is $\sim$18~MHz of centroid shift). \denis{}
fits this calibration itself on every load, so the numbers do not
depend on which \code{clstools} build is installed, and by default
\emph{nothing is dropped}: a run whose calibration outliers would
actually move its centroid gets a blinking \app{!} badge on its row
(in both Pre-Analysis and Analysis). Clicking it reports the outlier
count, the centroid shift in MHz, and the axis tilt across the scan,
with \app{Open calibration\dots} / \app{Keep as-is} --- either way
the blinking stops for good, and the acknowledgement is saved with
the project.

Right-clicking a run and choosing \app{Calibration\dots} opens a dialog
for the DAC$\rightarrow$HV voltage calibration in force for that run. By
default the app fits every calibration point itself; the override is a
property of the file and is stored in the project. Under
\app{Fit this run's own points} you may raise the polynomial order,
exclude points (drop the first $N$ settling points, an $n\!\cdot\!\sigma$
residual cut, or an explicit hand-picked set), or tick
\app{Force zero intercept (p\textsubscript{0} = 0)} to constrain the
calibration offset to zero for a chain known to have no DC offset --- the
offset is dropped from the fit, not fitted and then blanked, so the
gain is the true through-origin slope. You can instead borrow another
run's calibration or type the coefficients in directly. The dialog shows
the resulting polynomial, the residuals, and the centroid shift the
choice costs in MHz. \app{Calibration overview\dots} triages every
loaded run at once, worst first. Overrides live in a process-wide
registry shared by every tab and persist under the save file's
top-level \code{calibrations:} key (acknowledgements under
\code{calibration\_acks:}), siblings of \code{scan\_filters:}.
